\begin{table}[H]
\scriptsize
\setlength{\tabcolsep}{2.2pt}
\renewcommand{\arraystretch}{0.92}

\caption{\label{tab:ea20-policy-category-contributions}Annual effects of household price-policy categories on average inflation and inflation inequality in the euro area, 2021--2023}
\centering
\resizebox{\textwidth}{!}{%
\begin{tabular}[t]{l>{\centering\arraybackslash}p{1.6cm}*{6}{>{\centering\arraybackslash}p{1.15cm}}}
\toprule
\multicolumn{2}{c}{} & \multicolumn{3}{c}{\shortstack{Effect of price policies on\\average inflation}} & \multicolumn{3}{c}{\shortstack{Effect of price policies on\\Q1--Q5 inflation inequality}}\\
\cmidrule(lr){3-5} \cmidrule(lr){6-8}
Policy category & \shortstack{Fiscal cost\\(\% of\\2022 GDP)} & 2021 & 2022 & 2023 & 2021 & 2022 & 2023\\
\midrule
\multicolumn{8}{l}{\textit{Instrument classification}}\\
VAT and non-fuel excise cuts & -- & -0.1 & -0.4 & -0.2 & 0.0 & -0.2 & -0.1\\
Administered prices & -- & 0.0 & -0.7 & -0.7 & 0.0 & -0.4 & -0.4\\
Fuel price reductions & -- & 0.0 & -0.3 & 0.3 & 0.0 & 0.1 & 0.0\\
Other calibrated layer & -- & 0.0 & -0.1 & 0.0 & 0.0 & 0.0 & 0.0\\
\textbf{Total} & \textbf{1.58} & \textbf{-0.1} & \textbf{-1.5} & \textbf{-0.6} & \textbf{-0.1} & \textbf{-0.6} & \textbf{-0.5}\\
\midrule
\multicolumn{8}{l}{\textit{Targeted vs Non-targeted}}\\
Untargeted price measures & -- & -0.1 & -1.0 & -0.4 & 0.0 & -0.3 & -0.5\\
Targeted price measures & -- & 0.0 & -0.5 & -0.2 & 0.0 & -0.3 & -0.1\\
\textbf{Total} & \textbf{1.58} & \textbf{-0.1} & \textbf{-1.5} & \textbf{-0.6} & \textbf{-0.1} & \textbf{-0.6} & \textbf{-0.5}\\
\bottomrule
\end{tabular}
}%
\begin{flushleft}
\footnotesize{Notes. Observed minus counterfactual annual inflation or Q1--Q5 gap, in percentage points; negative values indicate a reduction. Fuel price reductions include all road-fuel (CP0722) relief, including excise-duty reductions, regulated-price support, and direct subsidies; the VAT and non-fuel excise category excludes road-fuel measures. Category effects need not sum to Total because policies interact. Targeting classifies measures, not individual eligibility. Fiscal costs are recorded policy envelopes divided by nominal 2022 GDP for the same 20 countries; budget periods may span several years. Fiscal costs refer to entire schemes affecting household prices, not amounts received exclusively by households. A dash indicates an unidentified comparable cost, not zero. Sources: Eurostat, national sources, authors' calculations.}
\end{flushleft}
\end{table}
